\section{Exercise 2}
\textit{Let $F$ be a vector field on $\mathbb{R}^3$ given by:}
\[F(x,y,z) = (3x,2y,-z)\]
\begin{itemize}
    \item \textit{Calculate $\nabla \cdot F$ and $\nabla \times F$ and find all function $\phi$ such that $\nabla\phi = F$}
    \item \textit{Calculate the flux of F out of the unit sphere $S^2 = \{(x,y,z) \in \mathbb{R}^3 | x^2 + y^2 + z^2 = 1\}$}
\end{itemize}

The product of $\nabla$ and F, which is vector field, yields a scalar field, this can be found as:
\[\nabla \cdot \mathbf{F} = 
\frac{\partial F_x}{\partial x}(3x) +
\frac{\partial F_y}{\partial y}(2y) +
\frac{\partial F_z}{\partial z}(-z) = 3 + 2 - 1 = 4\]
This is also known as the divergence, this tells us something about the flux of the vector field, leaving a small
closed surface.
The curl of the vector field tells us about local rotation and can calculated as follows:
\[\nabla \times \mathbf{F} = 
\begin{pmatrix}
    \frac{\partial F_z}{\partial y} - \frac{\partial F_y}{\partial z}\\
    \frac{\partial F_x }{\partial z} - \frac{\partial F_z}{\partial x}\\
    \frac{\partial F_y}{\partial x} - \frac{\partial F_x}{\partial y}
\end{pmatrix}
=
\begin{pmatrix}
    \frac{\partial (-z)}{\partial y} - \frac{\partial (2y)}{\partial z}\\
    \frac{\partial (3x) }{\partial z} - \frac{\partial (-z)}{\partial x}\\
    \frac{\partial (2y)}{\partial x} - \frac{\partial (3x)}{\partial y}
\end{pmatrix}
= 
(0 - 0, 0 - 0, 0 - 0) = 0
\]
This tells us that the vector field is \textit{Irrotational} and also convervative. This also means that
the field is defined everywhere in the space.  
%theorem 3.1
If a vector field $\mathbf{F}$ is convervative, then the corresponding scalar field $\phi$ is
called the \text{potential} for $\textbf{F}$. 
The gradient consists of three partial derivates that must yield the vector field itself:
\[
\frac{\partial \phi}{\partial x} = 3x\\
\qquad
\frac{\partial \phi}{\partial y} = 2y\\
\qquad
\frac{\partial \phi}{\partial z} = -z\\
\]
Integrate the first term with respect to x, however we need to introduce arbitrary functions g(y,z) and h(z) during each integration, to recover any
lost information, as we cannot be sure if the origianl fuction contained any terms relating to y and z.
Because:
\[\nabla\phi = 
\begin{pmatrix}
\frac{\partial \phi}{\partial x}\\    
\frac{\partial \phi}{\partial y}\\
\frac{\partial \phi}{\partial z}    
\end{pmatrix},
\qquad
\textit{and}
\qquad
\nabla\phi = \mathbf{F}
\qquad
\Rightarrow
\qquad
\mathbf{F} =
\begin{pmatrix}
\frac{\partial \phi}{\partial x}\\    
\frac{\partial \phi}{\partial y}\\
\frac{\partial \phi}{\partial z}    
\end{pmatrix}
\]
This means that we must find some scalar function $\phi$, that when partial differentiated yields or vector field.
As listed above, the three partials, we can start by integration the x component. 
\[\int 3x \ dx \Rightarrow \phi = \frac{3x^2}{2} + g(y,z)\]
If we differentiate $\phi$ for y, the x term goes away and we're left with:
\[\frac{\partial \phi}{\partial y} = \frac{\partial}{\partial y} \left(\frac{3}{2}x^2 +g(y,z)\right) \Rightarrow \frac{\partial \phi}{\partial y} = \frac{\partial g(y,z)}{\partial y} = 2y \]
Because from our definition from ealier, $\frac{\partial \phi}{\partial y} = 2y$. This can be integrated:
\[\int g(y,z) \ dy = 2\frac{1}{2}y^2 + h(z) = y^2 + h(z) \Rightarrow \phi = \frac{3x^2}{2} + y^2 + h(z)\]
Same process as before, if we differentiate $\phi$ for z, the x and y terms goes away, and we're left with:
\[\frac{\partial \phi}{\partial z} = \frac{\partial}{\partial z} \left(\frac{3}{2}x^2 + y^2 + h(z)\right) \Rightarrow \frac{\partial \phi}{\partial z} = \frac{\partial h(z)}{\partial z} = -z\]
Again, from our ealier definition of the partial derivative of z.
Integrating this, becomes:
\[\int h(z) \ dz = -1\frac{1}{2}z^2  = -\frac{z^2}{2} + C\]
And combining all the terms, yields the final expression:
\[\phi = \frac{3x^2}{2}  + y^2 -\frac{z^2}{2} + C \]
\[\]
\\
To find the flux flowing out of the unit sphere, the divergence theorem can be used. This relates the flux throught the surface
to the volume integral of the divergence. Here we transform the surface S into a solid sphere or ball \textit{B}, but this also requires a change to spherical coordinates:
$(r, \theta, \phi)$
\[\iint_{S^2} \mathbf{F} \cdot \mathbf{n} \ dS = \iiint_B \nabla \cdot \mathbf{F} \ dV\]
from ealier, the divergence was calculated:
\[\nabla \cdot \mathbf{F} = 3 + 2 - 1 = 4\]
This yields the integral:
\[\iiint_B 4 \ dV = \int_{0}^{1}\int_{0}^{\pi}\int_{0}^{2\pi} 4r^2 sin(\phi) \ d\theta \ d\phi \ dr\]
Where $dV = r^2 sin(\phi) \ d\theta \ d\phi \ dr$, where the bounds of integration are defined as:
\[
0 \le r \le 1,
\qquad
0 \le \phi \le \pi,
\qquad
0 \le \theta \le 2\pi
\]
% \[\Phi = \int_{0}^{1}\int_{0}^{\pi}\int_{0}^{2\pi} 4r^2 sin(\phi) \ d\theta \ d\phi \ dr\]
First integral:
\[= \int_{0}^{2\pi} 4r^2 sin(\phi) \ d\theta \Rightarrow 4r^2 sin(\phi) \ \int_{0}^{2\pi} 1 \ d\theta \Rightarrow 4r^2sin(\phi) \left[\theta\right]_0^{2\pi} = 4r^2sin(\phi)(2\pi - 0) = 8\pi r^2 sin(\phi)\]
Second integral:
\[= \int_{0}^{\pi} 8\pi r^2 sin(\phi) \  d\phi \Rightarrow 8\pi r^2 \int_{0}^{\pi} sin(\phi) \ d\phi \Rightarrow 8\pi r^2 \left[-cos\right]_0^\pi\] 
\[\Rightarrow 8\pi r^2 ((-cos(\pi))-(-cos(0))) = 8\pi r^2 ((-(-1))-(-1)) = 16\pi r^2 \] 
Last and outer Integral
\[= \int_{0}^{1}16\pi r^2 \ dr \Rightarrow 16\pi \int_{0}^{1} r^2 \ dr \Rightarrow \left[\frac{16\pi}{3}r^3\right]_0^1 = \frac{16\pi}{3}\]
\[\]
\[\]

%https://www.clinicalbiostats.com/tutorials/undergraduate-mathematics-calculus-iii/flux-through-spheres-using-the-divergence-theorem/